\documentclass[a4paper,11pt]{ltjsarticle}
\usepackage[noto-otf]{luatexja-preset}
\usepackage{amsmath,amssymb,amsthm,mathtools}
\usepackage{booktabs,array,longtable,enumitem}
\usepackage[hidelinks]{hyperref}
\usepackage{url}
\usepackage[margin=25mm]{geometry}

\theoremstyle{plain}
\newtheorem{theorem}{定理}[section]
\newtheorem{proposition}[theorem]{命題}
\newtheorem{lemma}[theorem]{補題}
\newtheorem{corollary}[theorem]{系}
\theoremstyle{definition}
\newtheorem{definition}[theorem]{定義}
\newtheorem{remark}[theorem]{注意}
\newtheorem{convention}[theorem]{規約}
\renewcommand{\proofname}{証明}

\newcommand{\Lq}{\ulcorner\mathcal{L}\urcorner}
\newcommand{\cq}[1]{\ulcorner #1\urcorner}
\newcommand{\Lang}{\mathcal{L}}
\newcommand{\M}{\mathcal{M}}
\newcommand{\LNGN}{\mathrm{LNGN}}
\newcommand{\Con}{\mathrm{Con}}
\newcommand{\rank}{\operatorname{rank}}
\newcommand{\CNF}{\mathrm{CNF}}
\newcommand{\Def}{\mathrm{Def}}
\newcommand{\Num}{\mathrm{Num}}
\newcommand{\Bg}{\mathrm{Big}}
\newcommand{\ZFC}{\mathsf{ZFC}}
\newcommand{\ZF}{\mathsf{ZF}}
\newcommand{\PA}{\mathsf{PA}}
\newcommand{\Sth}{\mathsf{S}}
\newcommand{\T}{T'_{\LNGN}}
\newcommand{\code}[1]{\texttt{#1}}
\newcommand{\file}[1]{\texttt{#1}}

\title{月絶天葬数\\[0.5em]{\large Mahlo 基数で選んだ評価構造上の Rayo 型関数と，巨大数庭園数との比較}}
\author{葵}
\date{2026 年 10 月（審議用文書 改訂 68 にもとづく）}

\begin{document}
\maketitle

\begin{abstract}
構成可能宇宙 $L$ における最小の Mahlo 基数 $\kappa$ を用い，評価構造 $\M=(V_{\kappa+1})^L$ 上の純粋な $\{\in,=\}$ 一階式による一意定義可能性から Rayo 型の関数 $F$ を定め，自然数 $N=F^{F(10^{100})}(10^{100})$（月絶天葬数）を定義する．
定義に追加の記号・パラメータ・選択関数は用いない．
次に，巨大数庭園数（LNGN）を同じ宇宙 $L$ で評価し，その Henkin 公理を自由変数について修復した理論 $\T$ と，各宇宙 $U(\alpha)$ を $\alpha$ 番目の到達不能基数の段に置く候補解釈 $I$ を定める．
$I$ の下での値 $\LNGN_I$ を名指す長さ $19{,}229{,}718{,}256$ の式 $\varphi$ を具体的に構成し，$\M\models\exists!x\,\varphi(x)$ かつ $\M\models\varphi[\LNGN_I]$（定理 G）を示すことで，$\LNGN_I<F(10^{100})<N$ を得る．
定理 G の証明は，Lean 4 による部分形式化（式の書き写し・意味・一意性）と，$\ZFC$ における集合論的な存在証明（階数の評価と再帰の非循環性）を結合したものである．
あわせて，$\Sth=\ZFC+V{=}L+\text{「Mahlo 基数が存在する」}$ が $\Con(\T)$ を証明すること，および $V_{\lambda'+1}$ に対する $F$ の類似物を $F$ が具体的なしきい値から先で上回ることを示す．
\end{abstract}

\tableofcontents

\section{序論}

\subsection{背景}
Rayo の数 \cite{rayo} は，集合論の言語の式による自然数の定義可能性を用いて巨大数を定める．
すなわち，$k$ 記号未満の式で一意に定義できる自然数の最大値に $1$ を足す関数を考え，それに大きな引数を与える．
この種の定義では，「定義可能」を評価する構造の選び方が数の大きさを決める．
本稿では，評価構造そのものを大きな基数で選ぶ．
構成可能宇宙 $L$ における最小の Mahlo 基数 $\kappa$ をとり，$\M=(V_{\kappa+1})^L$ を評価構造に固定する．

比較の相手は巨大数庭園数（Large Number Garden Number，以下 $\LNGN$）\cite{lngn} である．
$\LNGN$ は，順序数 $\alpha$ ごとの宇宙 $U(\alpha)$ の列と Henkin 型の選択関数 $\Theta$ をもつ理論 $T$ を立て，各 $U(\mathrm{CNF}(i))$ での定義可能性から関数 $f$ を定め，$f^{10}(10\uparrow^{10}10)$ とする．

\subsection{主結果}
本稿の結果は次の三つである．

\begin{enumerate}[label=(\arabic*)]
\item \textbf{定義．}（第 \ref{sec:prelim}・\ref{sec:def} 節）$V$ に Mahlo 基数が存在するという前提の下で，$\M$ と関数 $F$ は $L$ の中で一意に定まり，$F$ は全域かつ広義単調増加で，$y\ge 10^{100}$ なら $F(y)>y$ である．よって
\[ N = F^{F(10^{100})}(10^{100}) \]
は自然数として確定する（月絶天葬数，$N_{\mathrm{Moon}}=N$）．
\item \textbf{条件つきの支配．}（第 \ref{sec:lngn}〜\ref{sec:proofG} 節）$\LNGN$ を $L$ で評価し，修復した理論 $\T$ の候補解釈 $I$ を固定するとき，
\[ \LNGN_I < F(10^{100}) < N. \]
これは，$\LNGN_I$ を名指す長さ $19{,}229{,}718{,}256<10^{100}$ の式 $\varphi$ を構成し，定理 G（$\M\models\exists!x\,\varphi(x)$ かつ $\M\models\varphi[\LNGN_I]$）を示すことによる．
\item \textbf{新規性の補強．}（第 \ref{sec:novel} 節）$\Sth\vdash\Con(\T)$ であり，$\Sth$ が無矛盾なら $\T\nvdash\Con(\Sth)$．また，$\M$ の中で短く名指せる $\lambda'<\kappa$ について，$V_{\lambda'+1}$ に対する $F$ の類似物 $F_{\lambda'}$ は，$k\ge 2(|\psi_{\lambda'}|+32{,}574{,}917)$ で $F_{\lambda'}(k)<F(k)$ を満たす．
\end{enumerate}

(2) は「比較の前提」（$\LNGN$ も $L$ で評価すること，解釈 $I$ を固定すること）の下での主張である．
$\LNGN$ は解釈ごとに値が変わり，許される解釈全体では値が上に有界でない（命題 \ref{prop:unbounded}）．
本稿は，$U(\alpha)$ を $\kappa$ より上に置く解釈や，解釈全体の上限については何も主張しない．

\subsection{証明の性格}
定理 G の証明は，次の二つを結合したものである．Lean による完全な形式化ではない．
\begin{itemize}
\item \textbf{Lean 4 による部分形式化．}生成した式 $\varphi$ の 117 個の部品が，集合の言葉で書いた意図した条件を正しく書き写していること．その条件の意味（Tarski の充足，定義可能性，$L_\lambda$ と $<_\lambda$，$\Theta$ の Henkin 公理）．$\M$ の中で $\varphi$ を満たす元が高々一つであること（一意性．仮定はモデルの公理だけ）．
\item \textbf{$\ZFC$ における紙の証明．}意図した証人（計算表・段階の表・構造の表）が $V_\kappa$ に属し，$\M$ の中で条件を満たすこと（存在）．階数の評価と，構造の再帰が循環しないことを含む．
\end{itemize}
Lean の定理は $L$ の中の $\ZFC$ の定理として読む（\S\ref{sec:readL}）．この読み替えはメタな議論で，機械では検査していない．
本稿では証明の一部を概略に圧縮した．詳細は審議用文書 \cite{doc} にあり，Lean のファイルと検査の記録はリポジトリ \cite{repo} にある．Lean のファイルは，本稿とは独立の環境（GitHub Actions）でも再コンパイルされた．$\varphi$ の全文（$57{,}684{,}403{,}816$ バイト）も，生成プログラムとは独立に書いた構文検査器で検査した（第 \ref{sec:verify} 節）．

\subsection{構成}
第 \ref{sec:prelim} 節で評価モデルを固定し，第 \ref{sec:def} 節で $F$ と $N$ を定義する．
第 \ref{sec:lngn} 節で $\LNGN$ を定式化し，第 \ref{sec:main} 節で主定理を述べる．
第 \ref{sec:proofG} 節で定理 G を証明し，第 \ref{sec:verify} 節で機械検証をまとめる．
第 \ref{sec:novel} 節は新規性の補強，第 \ref{sec:discussion} 節は比較の前提と未確定の点の議論である．

\section{準備：評価モデル}\label{sec:prelim}

\subsection{記法}
$V_\alpha$ は累積階層，$\rank(x)$ は階数で，$\rank(x)<\mu\iff x\in V_\mu$．
自然数はフォン・ノイマン自然数とし，$0=\emptyset$，$n+1=n\cup\{n\}$．
順序対は $(a,b)=\{\{a\},\{a,b\}\}$ とする．
\emph{1-到達不能基数}とは，到達不能基数であって，その下に到達不能基数が非有界に存在するものをいう．

\subsection{評価モデルの固定}

\begin{proposition}\label{prop:mahloL}
$V$ の Mahlo 基数 $\mu$ は $L$ でも Mahlo 基数である．
\end{proposition}
\begin{proof}
$V$ で正則な基数は $L$ でも正則である．$\mu$ は $V$ で極限基数で，$V$ の基数は $L$ でも基数なので，$\mu$ は $L$ でも極限基数である．$L$ では GCH が成り立つので，$\mu$ 未満の $L$ の基数 $\lambda$ について $(2^\lambda)^L=(\lambda^+)^L<\mu$ となり，$\mu$ は $L$ で強極限である．よって $\mu$ は $L$ で到達不能である．
$\mu$ 未満の $V$ の到達不能基数 $\nu$ は，同じ理由で $L$ でも到達不能である．$L$ に属する $\mu$ の閉非有界集合は $V$ でも閉非有界であり，$\mu$ は $V$ で Mahlo なので，$V$ の到達不能基数の集合はそれと交わる．したがって $L$ の到達不能基数の集合は，$L$ のどの閉非有界集合とも交わり，$L$ で定常である．
\end{proof}

したがって，$V$ に Mahlo 基数が存在すれば $L\models\Sth$ である．ここで
\[ \Sth = \ZFC + V{=}L + \text{「Mahlo 基数が存在する」}. \]
以後，すべての定義と評価は $L$ の中で行う．

\begin{definition}
$\kappa$ を $L$ における最小の Mahlo 基数とし，
\[ \M = (V_{\kappa+1},\in) \quad (L \text{ の中で計算する}) \]
とおく．$V{=}L$ の下で $V_\kappa=L_\kappa$ である．
\end{definition}

$\omega$ は $L$ と $V$ で一致するので，以下で定まる自然数は $V$ の自然数としても同じものである．
本稿の主張は「$L$ で評価する」という意味論上の取り決めの下でのものである．これは，Rayo の数が真の宇宙 $V$ を基準とするのと同じ種類の取り決めである．

\section{月絶天葬数の定義}\label{sec:def}

\subsection{言語と記号数}
基本記号のアルファベットを $\Sigma=\{\,x,\ {}',\ \neg,\ \wedge,\ \forall,\ \in,\ =,\ (,\ )\,\}$ とする．
変数は $x$ のあとに有限個のプライムを並べた語 $x, x', x'', \dots$ である．
\emph{$\Lang$-式}は次の規則で生成されるものに限る．
\[ \varphi ::= (u\in v) \mid (u=v) \mid \neg\varphi \mid (\varphi\wedge\psi) \mid \forall u\,\varphi \]
$\exists,\vee,\to,\leftrightarrow$ は略記で，記号数は展開後の式について数える．
式 $\varphi$ の\emph{記号数} $|\varphi|$ は，変数を構成する $x$ とプライムを含め，基本記号を一つずつ数えた総数である．

\subsection{一意定義可能性と関数 $F$}
\begin{definition}
自然数 $n$ が\emph{$k$ 記号未満の式で $\M$ において一意に定まる}とは，自由変数がちょうど $x$ である $\Lang$-式 $\varphi(x)$ で
\[ |\varphi|<k,\qquad \M\models\exists!x\,\varphi(x),\qquad \M\models\varphi[n] \]
を満たすものが存在することをいう．充足関係は，$\M$ が集合であることを用いて $L$ の中で Tarski の再帰で定める．数 $n$ は対象言語の項として書かず，割当て $x\mapsto n$ で指定する．
\end{definition}

\begin{definition}
$D_k=\{\,n\in\omega : n$ は $k$ 記号未満の式で $\M$ において一意に定まる$\}$ とし，
\[ F(k) = \begin{cases} 1+\max D_k & (D_k\neq\emptyset)\\ 1 & (D_k=\emptyset)\end{cases} \]
と定める．
\end{definition}

\begin{proposition}
$F$ は全域で，広義単調増加である．
\end{proposition}
\begin{proof}
長さ $k$ 未満の語は有限個しかなく，各式が一意に定める自然数は高々一つなので，$D_k$ は有限である．また $D_k\subseteq D_{k+1}$ である．
\end{proof}

Mahlo 基数の存在と $V{=}L$ は，$F$ の全域性のためではなく，評価構造 $\M$ を固定するために用いる．

\subsection{補題：$y\ge 10^{100}$ で $F(y)>y$}
$y$ そのものを名指す必要はなく，$y$ 以上の数を短く名指せば足りる．そのために $|V_j|$ を用いる．

\begin{definition}
次の部品から式 $\Bg_j(x)$ を組み立てる（$\exists,\vee,\to,\leftrightarrow$ は展開する）．
\begin{itemize}
\item $\mathrm{Empty}(a):\ \forall e\,\neg(e\in a)$．$\mathrm{Trans}(a):\ \forall z\,(z\in a\to\forall u\,(u\in z\to u\in a))$．
\item $\mathrm{Succ}(z):\ \exists w\,(w\in z\wedge\forall u\,(u\in z\leftrightarrow(u\in w\vee u=w)))$．
\item $\mathrm{Nat}(a)$：$\mathrm{Trans}(a)$，$a$ の元も推移的，$a$ 以下の各元が $\mathrm{Empty}$ か $\mathrm{Succ}$．
\item $\mathrm{PowerOf}(a,b):\ \forall z\,(z\in a\leftrightarrow\forall u\,(u\in z\to u\in b))$．
\item $\mathrm{Chain}_0(v)=\mathrm{Empty}(v)$，$\mathrm{Chain}_{j+1}(v)=\exists w\,(\mathrm{PowerOf}(v,w)\wedge\mathrm{Chain}_j(w))$（$v,w$ には $x'$ と $x''$ を交互に使う）．
\item $\mathrm{Equi}(a,b)$：$a$ と $b$ の間の全単射が存在する．
\end{itemize}
$\Bg_j(x)\equiv \mathrm{Nat}(x)\wedge\exists v\,(\mathrm{Chain}_j(v)\wedge\mathrm{Equi}(x,v))$．
\end{definition}

\begin{lemma}\label{lem:big}
$\M\models\Bg_j[x]$ となる $x$ は $|V_j|$ だけである．また $|\Bg_j|\le 6659+114j$ で，$|\Bg_6|=7334$．
\end{lemma}
\begin{proof}
$\M$ は推移的で基礎の公理を満たす．推移的でその元も推移的な集合は順序数であり，$x$ 以下の順序数がすべて $0$ か後者なら $x$ は自然数である．$V_j$ は遺伝的有限で，その部分集合はすべて $\M$ に属するので，$\M$ の中の $\mathrm{PowerOf}$ は真の冪集合を与え，$\mathrm{Chain}_j(v)$ を満たす $v$ は $V_j$ だけである．$x$ と $V_j$ はともに有限なので，その間の全単射は $\M$ に属する．記号数はプログラム（\ref{app:repro}）で数えた．文字列の長さと構文木の再帰計算の二通りで数え，一致を確かめている．
\end{proof}

\begin{proposition}\label{prop:Fgrow}
$y\ge 10^{100}$ なら $F(y)>y$．
\end{proposition}
\begin{proof}
$|V_j|\ge y$ となる最小の $j$ をとる．$|V_5|=65536<y$ より $j\ge 6$．$j=6$ なら $|\Bg_6|=7334<y$．$j\ge 7$ なら $y>|V_{j-1}|$ である．$k\ge 6$ で $|V_k|\ge 2^{k+40}$ なので，$y>2^{j+39}$．一方 $|\Bg_j|\le 6659+114j<2^{13}(1+j)\le 2^{j+13}<2^{j+39}$．よって $|V_j|\in D_y$ かつ $|V_j|\ge y$ で，$F(y)\ge 1+|V_j|>y$．
\end{proof}

\subsection{月絶天葬数}
\begin{definition}
$N = F^{F(10^{100})}(10^{100})$．ここで $F^0(x)=x$，$F^{m+1}(x)=F(F^m(x))$ である．
\end{definition}
命題 \ref{prop:Fgrow} より，反復の各段の値は $10^{100}$ 以上で，$F$ は各段で真に増加する．$F(10^{100})\ge 2$ なので $F(10^{100})<N$ である．

\section{巨大数庭園数の定式化}\label{sec:lngn}

\subsection{原文の要約}
$\LNGN$ \cite{lngn} の理論は，$\in$ と一項関数記号 $U$ をもつ言語 $L$ 上の $T=\ZF_L\cup\{U1,U2,U3\}$ である．
各 $U(\alpha)$ は，内部の言語 $\Lq$ の構造である．$\Lq$ は $\cq{\in}$，関数記号 $\cq{U}$ と $\cq{\Theta}$，追加の定数・関数・関係記号 $\cq{c_i},\cq{f_{a,i}},\cq{r_{a,i}}$ をもつ．
\begin{itemize}
\item $U1$：任意の順序数 $\alpha$ について，$U(\alpha)\models\overline{\ZFC}_{\Lq}$．この理論は Henkin 公理 $\mathrm{Henkin}_{i,P}$ を含む．$P$ は $\Theta$ を含まない $\Lq$-式（Formula$_-$）で，公理の内容は「$P$ を満たす $x_i$ があれば $x_i:=\Theta(\mathrm{MetaCode}(P))$ でも $P$ が成り立つ」である．
\item $U2$：$\beta\in\alpha$ なら $\cq{U}^{U(\alpha)}(\beta)=U(\beta)$．
\item $U3$：ある $\beta$ について $|U(\alpha)|=V_\beta$ で，$\cq{\in}^{U(\alpha)}$ は $\in$ と一致する．
\end{itemize}
記号は $\omega^4$ 未満の順序数で符号化される（\ref{app:code}）．$\mathrm{MetaCode}(P)$ は記号の符号の列を $\mathrm{Encode}$ で自然数にしたもので，$\mathrm{CNF}:\omega\to\varepsilon_0$ はカントール標準形による全射である．
関数 $f$ と $\LNGN$ は次のように定められる．
\begin{itemize}
\item $\mathrm{IsDefinition}(P)$：$\exists x_0\,(P\wedge\forall x_1\,(P[x_1/x_0]\to x_1=x_0))$．
\item $\mathrm{Definable}(m,i,P)$：$i\in\omega$，$P$ は $\Lq$-式，$U(\CNF(i))\models\mathrm{IsDefinition}(P)$ かつ $U(\CNF(i))\models P[\mathrm{Parameter}(m)/x_0]$．
\item $f(n)=\sum\{\,m\in\omega : \exists i\in n\ \exists P\,(\mathrm{MetaCode}(P)\in n\wedge\mathrm{Definable}(m,i,P))\,\}$．
\item $\LNGN = f^{10}(10\uparrow^{10}10)$．
\end{itemize}

\subsection{well-definedness と修復}

\begin{proposition}[Henkin 公理と自由変数]\label{prop:henkin}
「構造が式を満たす」をすべての割当てで成り立つと読むと，$\overline{\ZFC}_{\Lq}$ はモデルをもたず，$U1$ により $T$ は矛盾する．
\end{proposition}
\begin{proof}
$i=1$，$P\equiv(\cq{x_1}\,\cq{\in}\,\cq{x_2})$ とする．$x_2$ に相異なる一元集合 $\{a\},\{b\}$ を割り当てると，$\mathrm{Henkin}_{1,P}$ は $\Theta(\mathrm{MetaCode}(P))=a$ と $=b$ を同時に要求する．
\end{proof}

この反例は $\Theta$ を含まないので，Formula$_-$ への制限の下でも成り立つ．そこで本稿は次の修復を採る．
\begin{convention}\label{conv:reading}
$\models$ は全称閉包で読む．Henkin 公理の対象は，自由変数が $\cq{x_i}$ だけの Formula$_-$ の式に限る（読み 1）．こうして修復した理論を $\T$ と書く．
\end{convention}

\begin{proposition}[Definable の一意性]
規約 \ref{conv:reading} の下で，同じ $i,P$ について $\mathrm{Definable}(m,i,P)$ を満たす $m$ は高々一つである．したがって $f(n)$ は有限和で，$f$ は well-defined である．
\end{proposition}
\begin{proof}
$U=U(\CNF(i))$ とし，$x_0,x_1$ 以外の自由変数への割当て $\tau$ を固定する．
全称閉包により，$P[\mathrm{Parameter}(m)/x_0]$ は $x_1$ への任意の値 $b$ と $\tau$ の下で真である．$b=m$ とすれば $P(x_0{=}m,x_1{=}m)$ が真である．
$\mathrm{IsDefinition}(P)$ の証人を $a$ とする．内側の $\forall x_1$ で $x_1=m$ とおくと，$P[x_1/x_0]$ は（原文の代入規則で束縛変数は付け替えられるので）$P(x_0{=}m,x_1{=}m)$ と同じ真理値をもち，真である．よって $m=a$ である．同様に $m'=a$ となる．
\end{proof}

「ある割当てで成り立つ」と読むと，$P\equiv(x_0=x_2)$ によりすべての $m$ が Definable になり，$f$ は定まらない．

\begin{proposition}[解釈全体にわたる非有界性]\label{prop:unbounded}
$\kappa$ が $L$ で Mahlo なら，$A=10\uparrow^{10}10$ 以上の任意の自然数 $m$ について，$\T$ のモデル $(V_\kappa,\in,U_{I_m})$ で $\LNGN_{I_m}\ge m$ となるものがある．
\end{proposition}
\begin{proof}
解釈 $I_m$ を，\S\ref{sec:candidate} の候補解釈 $I$ と同じ超限再帰で作る．違いは一点だけで，各 $U(\alpha)$ で追加の定数記号 $\cq{c_0}$ を $m$ と解釈する．$m\in V_\omega\subseteq V_{\lambda_\alpha}$ なので，これは構造である．

再帰をやり直すのは，次の二つのためである．
\begin{itemize}
\item $\Theta$ は再計算する．$\cq{c_0}$ を含む Formula$_-$ の式があるので，各段の $\Theta(m')$ は，$\cq{c_0}\mapsto m$ とした構造で good な元のうち $<_L$ で最小のものとして取り直す．たとえば $P(x_0)\equiv(x_0\,\cq{=}\,\cq{c_0})$ では $\Theta(\mathrm{MetaCode}(P))=m$ となり，Henkin 公理はこの $P$ についても成り立つ．$\cq{c_0}$ の解釈を変えて $\Theta$ を元のままにすると，Henkin 公理が破れる．
\item $\cq{U}$ の値も作り直す．$\cq{U}^{U(\alpha)}(\beta)$ は，新しい再帰で作った $U(\beta)$ そのものである．
\end{itemize}
命題 \ref{prop:model} の証明は，追加の定数記号の解釈に依存しない．$U1$ の $\ZFC$ の図式は記号の解釈によらずに成り立ち，Henkin 公理は各段の $\Theta$ の作り方から従い，$U2$・$U3$ は再帰の定義から従う．したがって，$(V_\kappa,\in,U_{I_m})\models\T$ である．

$P(x_0)\equiv(x_0\,\cq{=}\,\cq{c_0})$ は，$U(0)=U(\CNF(0))$ で $\mathrm{IsDefinition}(P)$ を満たし，$m$ を定義する．よって $\mathrm{Definable}(m,0,P)$ である．
$\mathrm{MetaCode}(P)\le 10^{10^{23.82}}$ は $A$ よりはるかに小さい．したがって，$n>\mathrm{MetaCode}(P)$ を満たす $n$ について，$f(n)$ の和に $m$ が含まれ，$f(n)\ge m$ となる．
$f^0(A)=A>\mathrm{MetaCode}(P)$ から始めると，各段で $f^{j+1}(A)\ge m\ge A>\mathrm{MetaCode}(P)$ となる．よって $\LNGN_{I_m}=f^{10}(A)\ge m$ である．
\end{proof}

したがって $\LNGN$ を単一の自然数として扱うには，解釈を固定する必要がある．解釈 $I$ の下での値を $\LNGN_I$ と書く．

\subsection{候補解釈 $I$}\label{sec:candidate}
\begin{definition}[候補解釈]\label{def:I}
$\lambda_\alpha$ を $\alpha$ 番目の到達不能基数とする．構造 $U(\alpha)$ を次で定める．
\begin{itemize}
\item 領域 $D=V_{\lambda_\alpha}$，$\cq{\in}$ の解釈 $E=\in\restriction D$．
\item 追加の定数記号は $\emptyset$，関数記号 $\cq{f_{a,i}}$ は値がつねに $\emptyset$，関係記号は空．
\item $\cq{U}(x)$ は，$x$ が順序数で $x\in\alpha$ なら構造 $U(x)$ そのもの，それ以外は $\emptyset$．
\item $\Theta(m)$ は，$m$ が Formula$_-$ の式 $P$ の番号なら，$P$ について \emph{good} な元のうち $<_L$ で最小のもの．good の定義は次のとおり．$P$ のどの変数 $x_i$ についても，「$P$ の真偽が $x_i$ の値だけで決まり，$P$ を満たす割当てがある」ならば，すべての割当て $\sigma$ で $\sigma[i:=a]\models P$ となる元 $a$．$m$ が番号でなければ $\Theta(m)=\emptyset$．
\end{itemize}
構造は五つ組 $((D,E),(Cn,(Fn,Rel)))$ で表す（構造の符号）．
\end{definition}
$V{=}L$ の下で $V_{\lambda_\alpha}=L_{\lambda_\alpha}$ であり，$<_L$ はその上の定義可能な整列順序である（\S\ref{sec:proofG} で具体化する）．
読み 1 の Henkin 公理の対象 $P$ は真偽が $x_i$ だけで決まるので，$P$ が定数でなければ $\Theta(\mathrm{MetaCode}(P))$ は $P$ を満たす．
$P$ は $\Theta$ を含まないので，$\Theta$ の定義は循環しない．

\begin{proposition}[候補解釈が $\T$ のモデルであること]\label{prop:model}
$\kappa$ が $L$ で Mahlo なら，$U_I$ を $\alpha<\kappa$ について上で定め，順序数でない引数では $\emptyset$ とするとき，$(V_\kappa,\in,U_I)\models\T$．
\end{proposition}
\begin{proof}[証明の概略]
$\ZF_L$ の個々の公理は，$\kappa$ の到達不能性と $V_\kappa$ の推移性から従う．
分出図式は，$A\in V_\kappa$ の部分集合が $V_\kappa$ に属することから従う．置換図式は $\kappa$ の正則性から従う．
$U1$ について：$\lambda_\alpha$ は到達不能なので，$V_{\lambda_\alpha}$ は追加記号をどう解釈しても $\ZFC$ の図式を満たす．Henkin 公理（読み 1）は上の注意による．
$U2$ は定義そのものである．$U(\beta)$ の階数は $\lambda_\beta+7$ 以下（\S\ref{sec:exist}）なので，$U(\beta)\in V_{\lambda_\alpha}$ となる．
$U3$ は $\beta=\lambda_\alpha$ である．
$\kappa$ は Mahlo なので，到達不能基数は $\kappa$ の下で定常で，順序型 $\kappa$ に並ぶ．よってすべての $\alpha<\kappa$ について $U(\alpha)$ が定まる．
\end{proof}

\begin{remark}[局所性]
$f$ は $\CNF(i)<\varepsilon_0$ の $U(\CNF(i))$ しか参照しない．$U2$ によって参照される $U(\beta)$ も $\beta<\varepsilon_0$ である．よって $\LNGN_I$ の値は $V_{\lambda_{\varepsilon_0}+1}$ だけで決まり，外側の宇宙の取り方に依存しない．
\end{remark}

\section{主定理}\label{sec:main}

\subsection{一般形}
\begin{proposition}[条件つき支配]\label{prop:B}
$\LNGN$ を $L$ で評価する．解釈 $I$ の初期部分 $I\restriction\varepsilon_0$ が $\M$ の中でパラメータなしに $\ell$ 記号の式で一意に定義できるとする．
$C$ を，$f$ と $f^{10}(10\uparrow^{10}10)$ の形式化，および $I\restriction\varepsilon_0$ の定義式との結合の費用を含めた定数とし，形式化で使う変数番号はすべて $C$ 未満とする．
$(\ell+1)(C+1)<10^{100}$ なら $\LNGN_I<F(10^{100})<N$．
\end{proposition}
\begin{proof}
$I\restriction\varepsilon_0$ の各成分は $V_\kappa$ に属し，その充足関係は $\M$ の中で定義できる．$I\restriction\varepsilon_0$ の定義式に $f$ の形式化を合成すれば，$\LNGN_I$ は $\M$ で一意に定まる．
合成の際，変数番号を $d\le C$ だけずらすと，長さは高々 $\ell(1+d)$ になる．したがって合成した式の長さは $\ell(C+1)+C<k=(\ell+1)(C+1)$ である．
$F$ は単調なので，$F(10^{100})\ge F(k)>\LNGN_I$ である．後半は第 \ref{sec:def} 節による．
\end{proof}

\subsection{式 $\varphi$ の具体的な構成}
本稿は，命題 \ref{prop:B} の $\ell$ と $C$ を経由せず，候補解釈 $I$ の下で $\LNGN_I$ を定める $\Lang$-式 $\varphi(x)$ そのものを構成した．
$\varphi$ は部品の合成で，部品はそれぞれ「条件を満たす有限な計算表・段階の表・構造の表が存在し，その値が\ldots」という形の式である．
主な部品と記号数を表 \ref{tab:parts} に示す．

\begin{table}[htbp]
\centering\small
\caption{$\varphi$ の主な部品と記号数}\label{tab:parts}
\begin{tabular}{p{0.27\textwidth}p{0.40\textwidth}r}
\toprule
部品 & 内容 & 記号数\\
\midrule
数詞 $A$ & $x=10\uparrow^{10}10=H(12,10,10)$（ハイパー演算の表） & 17{,}915\\
算術 & 加法・乗法・冪・素数・$n$ 番目の素数 & 例：乗法 23{,}461\\
Encode & $\mathrm{Encode}(a)=-1+2^h\prod p_{i+1}^{a_i}$ & 1{,}255{,}145\\
構文 & Formula・Formula$_-$（構文解析の証明書） & 10{,}119{,}792\\
充足関係 & 一般の $\Lq$ 構造での Tarski の充足 & 86{,}031{,}229\\
IsDefinition・Definable & 全称閉包・原文の代入規則 & 783{,}308{,}393\\
CNF の大小 & 名前どうしの比較の計算表 & 6{,}107{,}808\\
O と L & $O(i)=\CNF(i)$，$\lambda_\beta$ & 25{,}726{,}353 / 31{,}306{,}190\\
$V_\lambda$ & 冪集合と和集合の段階の表 & 19{,}550\\
$L_\lambda$ と $<_\lambda$ & 定義の表と段階の表 & 2{,}231{,}892{,}867\\
$\Theta$ & least\_good と theta\_L & 671{,}422{,}147\\
構造 $U(\CNF(i))$ & 構造の再帰表（struct\_full） & 7{,}594{,}788{,}010\\
$f$ & 構造の表・定義可能集合・有限和 & 19{,}104{,}908{,}587\\
\midrule
$\varphi$ 全体 & $f^{10}(A)$ & \textbf{19{,}229{,}718{,}256}\\
\bottomrule
\end{tabular}
\end{table}

$\varphi$ は第 3 節の文法の $\Lang$-式で，自由変数は $x$ だけである．
記号のほとんど（99.98\%）は変数のプライムである．生成プログラムが束縛変数に毎回新しい番号を割り当てるので，変数名が長くなる．

\subsection{主定理}
\begin{theorem}[定理 G]\label{thm:G}
$\M\models\exists!x\,\varphi(x)$ かつ $\M\models\varphi[\LNGN_I]$．
\end{theorem}

\begin{corollary}[主定理]\label{cor:main}
$\LNGN$ を $L$ で評価し，候補解釈 $I$（定義 \ref{def:I}，規約 \ref{conv:reading}）を固定するとき，
\[ \LNGN_I < F(10^{100}) < N. \]
\end{corollary}
\begin{proof}
$|\varphi|<10^{100}$ と定理 G より $\LNGN_I\in D_{10^{100}}$，よって $\LNGN_I<F(10^{100})$．後半は第 \ref{sec:def} 節による．
\end{proof}

\section{定理 G の証明}\label{sec:proofG}

以下はすべて $L$ の中の議論である．$H=V_\kappa$ と書く．

\subsection{$\M$ の性質と三つの原理}
\begin{enumerate}[label=(F\arabic*)]
\item $\M$ は推移的で $V_\kappa\subseteq\M$．$L$ に属する $V_\kappa$ の部分集合はすべて $\M$ に属する．
\item $L$ の中で $\lambda$ が到達不能なら $V_\lambda=L_\lambda$（$V{=}L$ の下で $H_\mu=L_\mu$，\cite[13 章]{jech}）．
\item $\M\subseteq V_{\kappa+1}$ なので，$\M$ の元の元はすべて $V_\kappa$ に属する．
\end{enumerate}

\paragraph{原理 A（$\Delta_0$ の絶対性）．}量化子がすべて有界な式は，推移的な $\M$ と $L$ の間で絶対的である．
順序対の判定 $\mathrm{pair}(p,a,b)$ は，「$p$ のどの元も $\{a\}$ か $\{a,b\}$ で，両方が $p$ に属する」という $\Delta_0$ の形で書く．\footnote{「$\forall s\,(s\in p\leftrightarrow s=\{a\}\vee s=\{a,b\})$」という有界でない形は，$a$ または $b$ の階数が $\kappa$ のとき偽陽性を生じる（たとえば $p=\emptyset$）．この点は審査で指摘され，改訂 29 で $\Delta_0$ 版に改めた．}

\paragraph{原理 B（有界でない量化子）．}部品の中の有界でない量化子は，次のどれかの形をしている．
(B0) 値を一意に固定する存在量化（$\emptyset$，数詞，後者，$\omega$）．
(B1) ガードつき全称量化で，ガードが「$W\in V_\kappa$」を含意するもの（割当て，有限関数，記号の 4 つ組，自然数，$\beta$ から $\lambda\le\kappa$ への関数）．
(B2) 同値型の全称量化 $\forall W\,(A(W)\leftrightarrow B(W))$ で，両辺が「$W\in V_\kappa$」を含意するもの．
(B3) 表や証人の存在量化．
(B0)〜(B2) では，量化子が $\M$ を動いても $L$ を動いても真偽は同じである．(B3) は原理 C で扱う．

\paragraph{原理 C（表の一意性と存在）．}部品が「条件 $\Psi(T,\vec a)$ を満たす表 $T$ が存在し，その値が\ldots」という形のとき，次の二つを示せば，$\M$ での真偽は意図した関係と一致する．
\begin{itemize}
\item \emph{一意性}：$\Psi$ を満たす $T\in\M$（階数 $\kappa$ の元を含む）の値は，意図した値に限る．
\item \emph{存在}：意図した表 $T^*$ が $V_\kappa$ に属し，$\M$ で $\Psi$ を満たす．
\end{itemize}

\subsection{部品の意味}
各部品の条件が意図した関係を表すことは，部品ごとの補題 G1〜G16 による（表 \ref{tab:lemmas}）．
これらの「条件 $\Leftrightarrow$ 意図した意味」の両方向は，ほぼすべて Lean で示した（第 \ref{sec:verify} 節）．

\begin{table}[htbp]
\centering\small
\caption{部品ごとの補題}\label{tab:lemmas}
\begin{tabular}{lp{0.55\textwidth}l}
\toprule
補題 & 内容 & 主な Lean の定理\\
\midrule
G1 & 空集合・後者・数詞・順序対・自然数・順序数の判定は，$\M$ の元について意図した関係 & \code{Export}\\
G2 & ハイパー演算・算術・素数・Encode の表 & \code{hyper\_iff} ほか\\
G3, G4 & 構文（Formula・Formula$_-$・MetaCode），一意読み & \code{formula\_iff}\\
G5 & 充足関係 $\Leftrightarrow$ Tarski の充足 & \code{sat\_iff}\\
G6 & IsDefinition・Definable（全称閉包，代入補題） & \code{definableIn\_iff}\\
G7 & CNF の大小の表 $\Leftrightarrow$ $\CNF(k)<\CNF(i)$ & \code{cnf\_cmp\_correct\_py}\\
G8〜G10 & 順序数，到達不能性，$O(i)=\CNF(i)$，$\lambda_\beta$ & \code{ordOf\_unique} ほか\\
G11 & $\mathrm{vstage}(\lambda,D)\Leftrightarrow D=V_\lambda$ & \code{vstage\_unique}\\
G12 & 構造の $\Theta$ 以外の成分 & \code{rep5\_struct}\\
G13 & $\mathrm{defines}$，$L_\lambda$ と $<_\lambda$ & \code{defines\_iff}，\code{lorder\_wo}\\
G14 & $\Theta$ は最小の good な元で Henkin 公理を満たす & \code{leastGood\_iff}\\
G15 & 構造 $U(\CNF(i))$ & \code{structFull\_unique}\\
G16 & $f$ と反復 & \code{fval\_unique}\\
\bottomrule
\end{tabular}
\end{table}

\subsection{一意性（Lean）}\label{sec:uniq}
\begin{theorem}[\code{full\_lngn\_unique'}，\file{TrUq4.lean}]\label{thm:uniq}
外延性・整礎性・自然数の数詞をもつ型 $\alpha$ 上の構造 $M$ が，閉包性 $\mathrm{Zax}$（$H$ の推移性・和・冪・分出・置換）と「$M$ の元の元は $H$ に属する」を満たすとする．Lean で組み立てた最終式 $\mathtt{full\_lngn}$ について，$M\models\mathtt{full\_lngn}(x)$ かつ $M\models\mathtt{full\_lngn}(x')$ なら $x=x'$．
\end{theorem}
証明は次の一意性の連鎖による．いずれも \code{sorry} を使わず，公理は Lean の標準の三つ（\code{propext}，\code{Classical.choice}，\code{Quot.sound}）だけである．
\begin{itemize}
\item $\omega$ の元はちょうど数詞である（\code{omega\_iff}）．
\item 集合の計算表は G7 の表の条件を満たす．よって $\mathrm{cnf\_lt}(k,i)\Rightarrow\CNF(k)<\CNF(i)$ である（\code{cnfLt\_lt}）．
\item $O$ の表は，名前の比較の整礎性による帰納法で一意である（\code{oOk\_unique}）．
\item $\varepsilon_0$，到達不能基数の列 $L$，$V_\lambda$ の表は，それぞれ一意である（\code{eps0\_unique}，\code{lamOk\_unique}，\code{vstage\_unique}）．
\item 構造：$\Theta$ 以外の関数記号の行は，下の段の表が一致すれば一致する（\code{fn\_agree}）．$\Theta$ の値は least\_good で決まる．Formula$_-$ の式は $\Theta$ を含まないので，式に現れる記号に限った構造の表現（\code{Rep5T}）を使うと，二つの $Fn$ の least\_good は同じ最小元になる（\code{theta\_agree}）．これを $\gamma$ に関する $\in$-帰納法と組み合わせて，\code{structRec\_unique} を得る．
\item $f$：構造の表と定義可能な集合は一意であり，有限和の値は上界の取り方によらない（\code{finiteSum\_val}）．反復の表は一意である（\code{iterate10\_unique}）．
\end{itemize}

\code{theta\_agree} は，構造の再帰が循環しないことの一意性の側の確認になっている．
$\Theta$ の値は，$Fn$ の $\Theta$ 以外の行だけで決まる．

\paragraph{閉包条件．}\code{rep5\_struct} は，$Fn$ の値がすべて $D$ に属することを使って $Fn$ から構造を作る．
このために必要な閉包条件（各段 $\gamma$ の $D$ は $\emptyset$ を含み，$x\in D\cap\gamma$ が順序数なら $SU(x)\in D$）も，Lean で証明した（\code{struct\_closed}）．
証明の材料は次のとおりである．
\begin{itemize}
\item 到達不能基数 $\lambda$ では $\beta\in\lambda\Rightarrow\beta+1\in\lambda$ である（\code{inacc\_lim}．強極限性と Cantor の対角線論法による）．
\item 極限の段 $V_\lambda$ は，数・順序対・有限列・記号の 4 つ組について閉じている（\code{vs\_*}）．
\item 二つの $V$ の段階の表は，共通の定義域で一致する（\code{vtab\_agree}）．
\item 以上を $\gamma$ に関する帰納法で組み合わせる．
\end{itemize}

\subsection{存在（紙の証明）}\label{sec:exist}
以下は $L$ の中の $\ZFC$ の議論である．意図した証人が $V_\kappa$ に属し，$\M$ で条件を満たすことを示す．
条件を $\M$ で読んでよいことは原理 A・B による．
この節では Lean の一意性の定理は使わない．使う Lean の定理は「条件 $\Leftrightarrow$ 意図した意味」の同値だけで，その「$\Leftarrow$」の向きを使う．

\paragraph{階数についての事実．}$\lambda$ が極限順序数なら次が成り立つ．
(R1) $a,b\in V_\lambda$ なら $(a,b)\in V_\lambda$．また，$\rank((a,b))=\max(\rank a,\rank b)+2$．
(R2) $V_\lambda$ の元の有限列は $V_\lambda$ に属する．
(R3) $V_\lambda$ の部分集合は $V_{\lambda+1}$ に属する．
(R4) 遺伝的有限な集合は $V_\omega$ に属する．

\begin{table}[htbp]
\centering\small
\caption{証人と階数の評価（$\lambda^*=\sup_{\beta<\varepsilon_0}\lambda_\beta<\kappa$）}\label{tab:rank}
\begin{tabular}{lp{0.43\textwidth}l}
\toprule
部品 & 証人 & 階数\\
\midrule
算術・ハイパー演算・CNF & 有限の計算表 & $<\omega$\\
$O$，$\varepsilon_0$ & $\{(i,\CNF(i)):i\in\omega\}$，$\varepsilon_0$ & $\le\varepsilon_0+2$\\
$L$ & $\{(\beta,\lambda_\beta):\beta<\varepsilon_0\}$ & $\le\lambda^*+2$\\
$V_\lambda$ & $\{(\beta,V_\beta):\beta\le\lambda\}$ & $\le\lambda+3$\\
$L_\lambda$，$<_\lambda$ & $\{(\beta,(L_\beta,<_\beta)):\beta\le\lambda\}$ & $\le\lambda+5$\\
定義の表 $DT_\beta$ & $\{((n,\tau),y)\}$ & $<\max(\beta,\omega)+\omega$\\
構造の表 & $SU^*=\{(\gamma,U(\gamma)):\gamma\le\alpha\}$ & $\le\lambda_\alpha+9$\\
$f$ の表 & 構造の表，有限集合 $S$，和と反復の表 & $<\lambda^*+\omega$\\
\bottomrule
\end{tabular}
\end{table}

$\lambda^*<\kappa$ である．$\kappa$ は Mahlo なので到達不能基数は $\kappa$ の下で非有界で，$\kappa$ は正則で $\varepsilon_0<\kappa$ だからである．
よって表 \ref{tab:rank} の証人はすべて $V_\kappa\subseteq\M$ に属する．
$DT_\beta$ の階数は，式の番号が自然数全体を動くので，$\beta=0$ でも $\omega$ 以上になる（たとえば $(x\in x),\neg\neg(x\in x),\dots$ はどれも $L_0=\emptyset$ 上で $\emptyset$ を定義する）．

\paragraph{構造の再帰と階数の同時帰納法．}
$U(\gamma)$（$\gamma\le\alpha$）を $\gamma$ に関する超限再帰で定める．同時に，次の主張を帰納法で示す．
\[ (\sharp_\gamma)\quad U(\gamma)=((D,E),(Cn,(Fn,Rel))),\ D=V_{\lambda_\gamma},\ \rank(U(\gamma))\le\lambda_\gamma+7. \]
各段 $\gamma$ では，$U\restriction\gamma$ と，すべての $\beta<\gamma$ についての $(\sharp_\beta)$ を仮定して，次の順に定める．
\begin{enumerate}
\item $\lambda=\lambda_\gamma$．$\beta<\gamma$ なら $\lambda_\beta<\lambda_\gamma$ である．
\item $D=V_\lambda$，$E=\in\restriction D$，$Cn=\{(c,\emptyset)\}$，$Rel=\emptyset$．
\item $Fn_0$ を，$f$ の行 $((f,w),\emptyset)$ と $U$ の行 $((U,\langle x\rangle),v_x)$ からなるものとする．$v_x$ は，$x\in D$ が順序数で $x\in\gamma$ なら $U(x)$，そうでなければ $\emptyset$ である．$(\sharp_x)$ より $\rank(U(x))\le\lambda_x+7<\lambda_\gamma$ なので，$U(x)\in D$ となる．
\item $W=<_\lambda$．
\item 構造 $\mathfrak A^0=(D,\in,\emptyset,Fn_0,\Theta\equiv\emptyset,\emptyset)$ で，各 Formula$_-$ の式について good な元のうち $W$ で最小のものを $v_m$ とする．Formula$_-$ の式の番号でない $m$ では $v_m=\emptyset$ とする．
\item $Fn=Fn_0\cup\{((\Theta,\langle m\rangle),v_m):m\in D\}$，$U(\gamma)=((D,E),(Cn,(Fn,Rel)))$．
\end{enumerate}
$(\sharp_\gamma)$ は次のように示せる．$E\subseteq V_\lambda$，$Cn\subseteq V_\omega$ である．$Fn$ の元 $((t,w),v)$ は $V_\lambda$ に属する（R1，R2，手順 3）．よって五つの成分は $V_{\lambda+1}$ に属し，組の階数は $\lambda+7$ 以下である．

\paragraph{非循環性．}
どの手順も，$U(\gamma)$ 自身や $U(\delta)$（$\delta\ge\gamma$）を参照しない．
$\Theta$ の値（手順 5）は $Fn_0$ だけから作った構造で決まり，$Fn$ の $\Theta$ の行は参照しない．
したがってこれは置換公理による通常の超限再帰で，$SU^*$ は集合として存在する．

\paragraph{存在の照合．}
意図した表が，生成した式の要求する条件を条項ごとに満たすことを確かめる．

\emph{$l\_order(\lambda,D,<_\lambda)$}：証人は $T^*=\{(\beta,(L_\beta,<_\beta)):\beta\le\lambda\}$ である．$\beta=0$，後者，極限の三つの場合を確かめる．
後者の段 $\beta=\beta_1+1$ では，定義の表 $DT^*=\{((n,\tau),y):\mathrm{def\_code}(n)\wedge\M\models\mathrm{defines}(n,\tau,L_{\beta_1},y)\}$ を証人にとる．
\begin{itemize}
\item def\_table の二条項は，$DT^*$ の取り方から直ちに従う．
\item 「defines で定義される集合の全体 $=\Def(L_{\beta_1})$」は，$\beta_1>0$ なら \code{defines\_iff}（$X\neq\emptyset$ を仮定する）と補題 \ref{lem:P} による．$\beta_1=0$ なら補題 \ref{lem:Z} による．
\item $<_{\beta}$ の三つの場合は，$<_L$ の定義そのものである．
\end{itemize}
$\lambda$ が到達不能なので，(F2) より $L_\lambda=V_\lambda=D$ である．

\emph{$\mathrm{struct\_rec}(\alpha,SU^*)$}：各段の level の条項を確かめる．
\begin{itemize}
\item lam\_of，vstage，e\_restrict，cn\_ok，$Rel=\emptyset$ は，それぞれ G10，G11 の存在の部分と定義から従う．
\item fn\_ok の五条項は，手順 3・5・6 から従う．$\Theta$ の部分は $\{(m,v_m)\}\subseteq D\times D$ による．関数性は，記号の種類が 4 つ組から一意に読めることによる．
\item theta\_L は次のように確かめる．
  \begin{itemize}
  \item $(\sharp_\gamma)$ より \code{rep5\_struct} の前提が満たされ，$Fn$ は構造 $\mathrm{sModel}(D,Fn)$ を表す．
  \item この構造と $\mathfrak A^0$ は，$\Theta$ 以外の記号の解釈が一致する．$\Theta$ を含まない式の充足は両者で一致するので，$v_m$ は least\_good を満たす（\code{leastGood\_iff} の「$\Leftarrow$」）．
  \end{itemize}
\end{itemize}
\code{rep5\_struct} の前提を満たすため，e\_restrict・cn\_ok・$Rel$・fn\_ok を theta\_L より先に確かめる．fn\_ok は $v_m\in D$ しか使わないので，この順序に循環はない．

\begin{lemma}[純粋な式の Def]\label{lem:P}
$X$ が推移的なら，等号を含む一階式で定義される $X$ の部分集合の全体と，記号が変数・$\in$・論理記号・括弧だけの式（純粋な式）で定義されるものの全体は一致する．
\end{lemma}
\begin{proof}
$u=v$ を $\forall z\,(z\in u\leftrightarrow z\in v)$ に置き換える．$X$ が推移的なら，$u,v\in X$ について外延性により同値である．束縛変数へのパラメータは充足に影響しない．
\end{proof}

\begin{lemma}[最初の段]\label{lem:Z}
$X=\emptyset$ のとき，$\mathrm{defines}(n,\tau,\emptyset,y)$ は「$n$ が $x_0$ 以外の変数を含まない純粋な式の番号で，$\tau=\emptyset$，$y=\emptyset$」と同値である．とくに定義される集合の全体は $\{\emptyset\}=L_1$ である．
\end{lemma}
\begin{proof}
条件を直接読む．割当て $\tau$ は空の定義域でしか存在しない．$y\subseteq\emptyset$ である．「$\forall z\in X$」は空虚に真である．$(x_0\in x_0)$ のような式は存在する．$<_1=\emptyset$ で，$\mathrm{new\_less}(\emptyset,\emptyset)$ は同じ定義どうしの比較なので偽である．
\end{proof}

\paragraph{Lean の定理の前提．}
存在の側で使う Lean の定理（\code{hyper\_iff}，\code{defines\_iff}，\code{lorder\_wo}，\code{rep5\_struct}，\code{leastGood\_iff}，\code{definableIn\_iff}，G14 の \code{good\_exists}）の前提を，一つずつ照合した．前提は次の三種類である．
\begin{itemize}
\item モデルの公理：Zax と hM．$L$ の中で $\M$ が満たす．
\item 集合が $H$ に属すること，推移性，空でないこと．
\item $s$ が式 $\varphi$ を表すこと（記号列の対応）．これは $\mathrm{metacode}$ と $\mathrm{formula}$ の補題から得る．
\end{itemize}
いずれも上の照合の中で満たされる．満たせない前提は，$X=\emptyset$ の段（補題 \ref{lem:Z} で別に扱った）のほかにはなかった．
$\mathrm{cnf\_lt}$ の表の存在は Lean になく，紙の議論による．G7.lean が計算する自然数の表を，遺伝的有限集合に写す．

\subsection{定理 G の結論}
$A=H(12,10,10)$ とし，反復の表 $G=\{(k,f^k(A)):k\le 10\}$ と，各 $f(n)$ の証人（構造の表，定義可能集合，和の表）をとる．
これらは $V_\kappa$ に属し，$\M$ で $\varphi[\LNGN_I]$ の条件を満たす．
定理 \ref{thm:uniq} と合わせて，$\M\models\exists!x\,\varphi(x)$ かつ $\M\models\varphi[\LNGN_I]$ である．\qed

\section{機械検証}\label{sec:verify}

\subsection{Lean による形式化の構成}
Lean 4（版 4.12.0，標準ライブラリのみ，Mathlib なし）\cite{lean} で 60 個のファイルを書き，カーネルで検査した．
\begin{itemize}
\item \textbf{抽象的な補題．}原理 A〜C と基本部品の絶対性（\file{Moon}，\file{Export}），充足関係（\file{G5}），定義可能性（\file{G6}），$L$ の段階の順序（\file{G13}），$\Theta$ と Henkin 公理（\file{G14}），CNF（\file{G7}，\file{Dec}），Encode と MetaCode（\file{Encode}）．
\item \textbf{書き写しの検証．}Python の生成プログラムの 117 部品を Lean の構文木として書き出し（\file{TrParts*}，\file{TrLit*}），各部品の意味が意図した条件と一致することを証明した（\file{TrRel*}，\file{TrMeans}）．展開の条件はカーネルの計算（\code{Eq.refl}）で検査する．
\item \textbf{意味への接続．}集合の言葉の条件から自然数の値（\file{TrNat}），構文（\file{TrSyn}），Tarski の充足（\file{TrSem}），定義可能性（\file{TrDef}，\file{TrL}），$L_\lambda$ と $<_\lambda$ の一意性と整列性（\file{TrLO*}），$\Theta$（\file{TrTh}），構造の表現（\file{TrU}）．
\item \textbf{一意性．}\file{TrUq1}〜\file{TrUq4}（\S\ref{sec:uniq}）．
\end{itemize}
\file{OldPair.lean} は，旧版の順序対で同じ証明が通らないことを示すためのファイルで，コンパイルが失敗するのが正しい．

\subsection{Lean の定理を読む宇宙}\label{sec:readL}
Lean の定理は「任意の型 $\alpha$ とその上の構造について」の形をしている．本稿はこれを $L$ の中で読み，$\alpha=\M$，$H=V_\kappa$ とする．
$L$ は $\ZFC$ のモデルで，$\kappa$ は $L$ で到達不能なので，$V_\kappa$ は任意の述語・関数について $\mathrm{Zax}$ を満たす．
$L$ には定義可能な大域的整列順序があるので，\code{Classical.choice} の解釈にも困らない．
この読み替え（Lean の論理の集合論的解釈）はメタな議論で，機械では検査していない．
なお，外側の $V$ で読むと，$V\neq L$ のとき分出・置換が成り立たない．たとえば $L$ に属さない実数 $r\subseteq\omega$ について，$\{n\in\omega:n\in r\}$ は $(V_\kappa)^L$ に属さない．

\subsection{最終式の照合}
\paragraph{ハッシュ値の照合．}$\varphi$ は木として展開すると約 192 億記号になる．
そこで，各節点のハッシュ値を節点の種類と子のハッシュ値から SHA-256 で作る Merkle 木の方式で，生成プログラムの出力と Lean で組み立てた式を比べた（\file{tr\_hash\_check.py}）．
結果は一致した．部品を一か所変えた変異版では不一致を検出した．

\paragraph{全文の独立な構文検査．}生成プログラムを使わず，第 3 節の文法だけから C で構文検査器を書いた（\file{phi\_syntax\_check.c}）．
$\varphi$ の全文を名前付きパイプで受け取り，一文字ずつ検査した．結果は次のとおりである．
\begin{itemize}
\item 全文はちょうど一つの $\Lang$-式である．
\item 記号数は $19{,}229{,}718{,}256$，自由変数は $\{x\}$ である．
\item 入れ子の最大の深さは 596，変数のプライムの最大個数は 44{,}481，全文は $57{,}684{,}403{,}816$ バイトである．
\end{itemize}

\subsection{独立の環境での再現}
付録一式を公開リポジトリ \cite{repo} に置いた．GitHub Actions が，GitHub の実行環境で次の二つを実行する．
\begin{itemize}
\item Lean 4.12.0 を新たに入れ，60 ファイルを順にコンパイルする．
\item ハッシュ照合と全文の構文検査を行う．
\end{itemize}
2026 年 10 月 21 日の実行（番号 39590487970）で，両方とも成功した．所要時間は Lean が 23 分，式の検査が 3 分である．

\subsection{検証の範囲}
機械で検査したのは，書き写し・意味・一意性と，最終式の構文・ハッシュの一致である．
存在の証明（\S\ref{sec:exist}）と，Lean を $L$ で読むこと（\S\ref{sec:readL}）は，人が読んで確かめる数学の証明である．
有限の総当たりと変異テスト（\ref{app:repro}）は，部品の書き間違いを見つけるための検査で，一般の構造での正しさの証明ではない．

\section{新規性の補強}\label{sec:novel}

この節の結果は，比較の規則上の「相手の数に使われていない理論」の根拠である．数の大小の根拠ではない．大小の根拠は，定義可能性による直接の対角化（第 \ref{sec:main} 節）だけに置く．

\subsection{原文の主張との関係}
$\LNGN$ の原文は，基礎理論の公理を強めても弱めても最終的な巨大数の大きさは変わらない，と述べる．本稿はこれに同意する．
本稿の新規性は，基礎理論に公理を足すことではない．Mahlo 基数 $\kappa$ によって，評価構造 $V_{\kappa+1}$ そのものを選ぶことにある．
構造が大きくなると名指せる集合が増える．たとえば $\LNGN$ の候補解釈全体も名指せるようになる．これが定義力の増大の源である．

\subsection{無矛盾性の強さ}
\begin{theorem}\label{thm:con}
$\Sth\vdash\Con(\T)$．さらに，$\Sth$ が無矛盾なら $\T\nvdash\Con(\Sth)$．
\end{theorem}
\begin{proof}
以下は $\Sth$ の中の議論である．
\begin{enumerate}[label=(\alph*)]
\item \emph{$\lambda_0$ の存在．}$\kappa$ を Mahlo 基数とする．$\kappa$ 未満の到達不能基数の集合は定常で，その極限点の集合は閉非有界である．両者は交わり，交わりの元は 1-到達不能基数である．最小のものを $\lambda_0<\kappa$ とする．
\item \emph{順序型．}$\lambda_0$ 未満の到達不能基数の集合は，正則な $\lambda_0$ で非有界である．よって順序型は $\lambda_0$ で，すべての $\alpha<\lambda_0$ について $\lambda_\alpha<\lambda_0$ となる．
\item \emph{モデル．}命題 \ref{prop:model} の証明が $\kappa$ について使うのは，到達不能性と「すべての $\alpha<\kappa$ で $\lambda_\alpha<\kappa$」だけである．よって $\mathcal N=(V_{\lambda_0},\in,U_I)\models\T$ となる．
\item \emph{絶対性．}$U1$ の「$U(\alpha)\models\overline{\ZFC}_{\Lq}'$」は，$\mathcal N$ の内部の充足関係で読まれる．$V_{\lambda_0}$ は推移的で $\omega$ を含む $\ZFC$ のモデルなので，集合の構造の充足関係は $V_{\lambda_0}$ と外側の宇宙の間で絶対的である．
\item \emph{一様性．}$\mathrm{Ax}_{\T}(n)$ を，$\T$ の公理の番号を表す $\Delta_1$ 式とする．示すべきことは
\[ \Sth\vdash\forall n\,(\mathrm{Ax}_{\T}(n)\to\forall s\,\mathrm{Sat}(\mathcal N,n,s)) \]
である．
\begin{itemize}
\item 個々の公理は，一つずつ示す．
\item 分出・置換の図式の例は，式の番号 $k$ を変数とする一つの式 $\mathrm{Sat}(\mathcal N,k,\cdot)$ に，$\Sth$ の図式を適用して一様に示す．
\item $U1$ は，(d) により「すべての $\alpha<\lambda_0$ と $\overline{\ZFC}_{\Lq}'$ の公理のすべての番号 $m$」についての主張に直して，同様に示す．
\end{itemize}
$\Sth$ の中で健全性定理を適用すると，$\Sth\vdash\Con(\T)$ を得る．
\item \emph{証明可能性．}$\T\vdash\Con(\Sth)$ とすると，$\PA\vdash\Con(\T)\to\Con(\Sth)$ である．理由は次のとおり．$\neg\Con(\Sth)$ は $\Sigma_1$ 文である．形式化した $\Sigma_1$ 完全性により $\PA\vdash\neg\Con(\Sth)\to\mathrm{Prov}_{\T}(\ulcorner\neg\Con(\Sth)\urcorner)$ となる．これと $\PA\vdash\mathrm{Prov}_{\T}(\ulcorner\Con(\Sth)\urcorner)$ を合わせればよい．
\item \emph{結論．}$\Sth$ は $\PA$ を（$\omega$ を通じた標準的な解釈で）証明する．よって $\Sth\vdash\Con(\T)\to\Con(\Sth)$ で，(e) と合わせて $\Sth\vdash\Con(\Sth)$ となる．これは第二不完全性定理に反する．この議論は $\T$ の無矛盾性を仮定しない．
\end{enumerate}
\end{proof}

\subsection{基準点に対する $F$ の類似物}
$\M$ の中で式 $\psi_{\lambda'}(\lambda)$ により一意に定まる $\lambda'<\kappa$ について，$F_{\lambda'}$ を $F$ の定義で $\M$ を $V_{\lambda'+1}$ に替えたものとする．

\begin{theorem}\label{thm:fsub}
$c_{\lambda'}=|\psi_{\lambda'}|+32{,}574{,}917$ とおく．$|\psi_{\lambda'}|$ は自由変数を $x''$ として数える．このとき，すべての $k\ge 1$ について
\[ F_{\lambda'}(k) < F(c_{\lambda'}+|\Num_k|),\qquad |\Num_k|\le 145+22{,}724\,(B-1) \]
が成り立つ．ここで $B$ は $k$ の二進の桁数である．とくに $k\ge 2(|\psi_{\lambda'}|+32{,}574{,}917)$ なら $F_{\lambda'}(k)<F(k)$．
\end{theorem}
\begin{proof}
「$W=V_{\lambda+1}$ について $x=F_W(k)$」を表す $\M$ の式 $G(\lambda,k,x)$ を構成した（\file{fsub.py}）．部品は次のとおりである．
\begin{itemize}
\item 第 3 節の文法の構文解析の証明書．
\item 変数の番号の表と自由変数の表．
\item $W$ 上の割当てについての充足の表．
\end{itemize}
証明書は一般には一意でない．根の構文解析に使わない真の項目を，余分に入れてよいからである．それでも次の三つが成り立つ．
\begin{itemize}
\item 証明書の項目はすべて真である．
\item 根の構文解析に現れる区間は，証明書に入っている．
\item 一意読みにより，自由変数と充足の表の値は，証明書の式の項目で真の値と一致する．
\end{itemize}
よって根の式についての判定は，表の取り方によらない．
$|G|=32{,}574{,}899$ である．$\Phi(x)\equiv\exists k\,(\Num_k(k)\wedge\exists\lambda\,(\psi_{\lambda'}(\lambda)\wedge G(\lambda,k,x)))$ のつなぎの記号は 17 個なので，$|\Phi|=|\Num_k|+|\psi_{\lambda'}|+32{,}574{,}916$ である．
$\M\models\exists!x\,\Phi(x)$ かつ $\M\models\Phi[F_{\lambda'}(k)]$ なので，$F_{\lambda'}(k)\in D_{|\Phi|+1}$ となる．
$\Num_k$ は Horner 法で組み立てる．一段は 21{,}625〜22{,}724 記号である．
$k\ge 2^{25}$ なら $145+22{,}724\log_2 k\le k/2$ なので，$k\ge 2c_{\lambda'}\,(\ge 2^{25})$ で $c_{\lambda'}+|\Num_k|\le k$ となる．$F$ は単調なので結論を得る．
\end{proof}
$G$ の正しさは，上の帰納法による紙の議論と有限の検証（\file{test\_fsub.py}，五つの層で 2{,}803 件）によるもので，Lean では検査していない．

\section{議論}\label{sec:discussion}

\subsection{比較の前提}
主定理（系 \ref{cor:main}）による比較は，次の二つが認められて初めて成り立つ．
\begin{enumerate}
\item \textbf{共通の評価宇宙．}$\LNGN$ も $L$ で評価する．$V\neq L$ のとき，$V$ で評価した $\LNGN$ は $L$ の外の集合を使いうるので，$\M$ の中では名指せない．
\item \textbf{解釈の固定．}$\LNGN$ の値を候補解釈 $I$ の下の値 $\LNGN_I$ に固定する．固定するものは次のとおりである．
  \begin{itemize}
  \item Henkin 公理の読み（規約 \ref{conv:reading}）．
  \item $U(\alpha)$ を $\alpha$ 番目の到達不能基数の段に置くこと．
  \item $\Theta$ と追加記号の解釈．
  \item 構造の符号（五つ組）．
  \end{itemize}
\end{enumerate}
公理 $U3$ は $|U(\alpha)|=V_\beta$ となる $\beta$ の存在しか要求しない．したがって，$U(\alpha)$ を $\kappa$ より上に置く解釈も許される．そのような解釈と，解釈全体の上限（命題 \ref{prop:unbounded} により有界でない）については，本稿は何も主張しない．

\subsection{原文の読みで確定していない点}
\begin{enumerate}
\item Henkin 公理を，$x_i$ 以外の自由変数をもつ式にどう適用するか（命題 \ref{prop:henkin}）．
\item 記号の符号化：4 つ組の長さ，等号の展開，関数・関係記号の適用の書き方（\ref{app:code}）．
\item $\mathrm{IsDefinition}(P)$ で，$P$ が $\cq{x_1}$ を自由にもつときに外側の $\forall\cq{x_1}$ が捕獲すること．$f$ の well-definedness には影響しないが，値には影響しうる．本稿は原文の代入規則どおりに実装した．
\end{enumerate}
なお，Okkuu の議論 \cite{okkuu} は $\Theta$ を含む式に Henkin 公理を適用している．現行の原文（対象は Formula$_-$）のままでは成立しない．命題 \ref{prop:henkin} の反例は $\Theta$ を含まないので，それとは独立である．

\subsection{前提の強さと $V{=}L$ の役割}
本稿の前提（$V$ に Mahlo 基数が存在する）は，無矛盾性の強さで $\LNGN$ の前提より真に強い（定理 \ref{thm:con}）．前提が強い分，成り立たない可能性も相対的に大きいことは認める．
$V{=}L$ は，評価モデルを $L$ に固定することと，$\Theta$ を $<_L$ で正準に定めることに使う．
評価モデルを固定しても，$\Con(\Sth)$ のような命題の独立性が解消されるわけではない．$N$ の値が一意なのは，「$L$ で評価する」という取り決めの下においてである．

\subsection{残る検査}
定理 G の証明のうち，存在の側（\S\ref{sec:exist}）と，Lean を $L$ で読むこと（\S\ref{sec:readL}）は，独立の読み手による検査を待つ．
機械で検査できる部分は，独立の環境での再コンパイルと全文の構文検査まで済んでいる．

\appendix
\section{$\Lq$ の記号と文法}\label{app:code}
原文は記号を $\omega^4$ 未満の順序数 $\omega^3a_3+\omega^2a_2+\omega a_1+a_0$ で符号化し，本稿はこれを 4 つ組 $(a_0,a_1,a_2,a_3)$ として扱う．
\begin{center}\small
\begin{tabular}{ll}
\toprule
記号 & 4 つ組\\
\midrule
$\bot,\to,\neg,\exists,(,)$ & $(0..5,0,0,0)$\\
$\cq{c_i}$ & $(6,0,i,0)$\\
$\cq{x_i}$ & $(7,1,i,0)$\\
$\cq{U}$，$\cq{\Theta}$ & $(8,2,0,1)$，$(8,3,0,1)$\\
$\cq{f_{a,i}}$ & $(8,4,i,a)$\\
$\cq{\in}$，$\cq{r_{a,i}}$ & $(9,5,0,2)$，$(9,6,i,a)$\\
\bottomrule
\end{tabular}
\end{center}
本稿は，原文の例 $(\cq{x_0}\cq{\in}\cq{x_1})$，$\cq{\Theta}\cq{(}\cq{x_0}\cq{)}$ に合わせて，文法を次のように読む．
\[ t ::= x_i\mid c_i\mid F\,(\,t_1\cdots t_a\,),\qquad \varphi ::= \bot\mid(t_1\in t_2)\mid R\,(\,t_1\cdots t_a\,)\mid(\varphi\to\psi)\mid\neg\varphi\mid\exists x_i\,(\varphi). \]
項は自己区切り的なので，コンマなしでも構文解析は一意である．
4 つ組の長さは常に 4（末尾の 0 を省かない）と読む．

\section{再現の手順}\label{app:repro}
検証用のプログラムと Lean のファイルは，リポジトリ \cite{repo} にある．主なものは次のとおりである．
\begin{center}\small
\begin{tabular}{p{0.38\textwidth}p{0.55\textwidth}}
\toprule
ファイル & 内容\\
\midrule
\file{num\_formula.py} & $\Bg_j$ の生成と記号数\\
\file{arith.py}，\file{parse.py}，\file{semantics.py}，\file{definable.py}，\file{cnf.py}，\file{ordinals.py}，\file{inacc\_seq.py}，\file{structure.py}，\file{lorder.py}，\file{ffunc.py} & $\varphi$ の部品の生成\\
\file{test\_*.py} & 有限の宇宙での総当たり検証と変異テスト\\
\file{final\_phi.py} & $|\varphi|$ と自由変数．\code{--stream} で全文を書き出す\\
\file{lean/*.lean}，\file{lean/build\_ci.sh} & Lean の形式化と，全ファイルのコンパイル\\
\file{tr\_hash\_check.py} & 最終式のハッシュ照合（一致なら終了コード 0）\\
\file{phi\_syntax\_check.c} & 全文の独立な構文検査\\
\file{fsub.py}，\file{test\_fsub.py} & 定理 \ref{thm:fsub} の式 $G$ と $\Num_k$\\
\file{.github/workflows/verify.yml} & 独立の環境での自動検査\\
\bottomrule
\end{tabular}
\end{center}
主な手順は次のとおりである．
\begin{itemize}
\item Lean の再コンパイル：\code{bash lean/build\_ci.sh}（Lean 4.12.0，約 25〜40 分）．
\item ハッシュ照合：\code{python3 tr\_hash\_check.py lngn}．
\item 全文の構文検査：\code{mkfifo phi.fifo; ./phi\_syntax\_check < phi.fifo \& python3 final\_phi.py --stream phi.fifo}（約 7 分．ディスクは不要）．
\end{itemize}


\section*{謝辞}
本稿の論証は，別の AI による複数回の独立の審査を受けて改められた．順序対の判定式の偽陽性（改訂 29），Henkin 公理の読みによる矛盾，$U2$ の実装の誤り（改訂 31），閉包条件の扱い，階数評価の誤り（改訂 55），証明書の一意性の誤った主張（改訂 65）などは，その審査の指摘による．

\begin{thebibliography}{9}
\bibitem{rayo} A.~Rayo, \emph{Big Number Duel}, Massachusetts Institute of Technology, 2007.
\bibitem{lngn} P進大好きbot，「高階集合論を超えた1階述語論理」，巨大数研究 Wiki ユーザーブログ，2019 年 12 月 3 日．\href{https://googology.fandom.com/ja/wiki/%E3%83%A6%E3%83%BC%E3%82%B6%E3%83%BC%E3%83%96%E3%83%AD%E3%82%B0:P%E9%80%B2%E5%A4%A7%E5%A5%BD%E3%81%8Dbot/%E9%AB%98%E9%9A%8E%E9%9B%86%E5%90%88%E8%AB%96%E3%82%92%E8%B6%85%E3%81%88%E3%81%9F1%E9%9A%8E%E8%BF%B0%E8%AA%9E%E8%AB%96%E7%90%86}{\nolinkurl{googology.fandom.com/ja/wiki/}ユーザーブログ:P進大好きbot/高階集合論を超えた1階述語論理}．巨大数庭園数の定義の原典（本稿は「理論」「埋め込み」「巨大関数」の各節を参照した）．英語の要約は Googology Wiki の記事 \url{https://googology.fandom.com/wiki/Large_Number_Garden_Number} にある．
\bibitem{okkuu} Okkuu，「巨大数庭園数の T と MK+ の矛盾について」（ブログ記事）．
\bibitem{jech} T.~Jech, \emph{Set Theory}, The Third Millennium Edition, Springer, 2003.
\bibitem{kanamori} A.~Kanamori, \emph{The Higher Infinite}, 2nd ed., Springer, 2003.
\bibitem{kunen} K.~Kunen, \emph{Set Theory}, College Publications, 2011.
\bibitem{lean} L.~de Moura and S.~Ullrich, The Lean 4 theorem prover and programming language, in: \emph{CADE-28}, LNCS 12699, Springer, 2021, pp.~625--635.
\bibitem{repo} 本稿の検証用リポジトリ，\url{https://github.com/ninnjinn008/moon-number-verify}．Lean のファイル・生成プログラム・検査器と CI の記録を含む（検証したコミットは 9dd73d5）．
\bibitem{doc} 葵，「巨大数庭園数に対抗する数 N の定義（審議用・改訂 68）」，\cite{repo} の \file{docs/} に収録．本稿で圧縮した証明（とくに定理 G の部品ごとの補題，存在の照合，Lean の定理の前提の照合）の詳細を含む．
\end{thebibliography}

\end{document}
